\begin{abstract}
As parameter-efficient fine-tuning (PEFT) experts proliferate, serving multiple
tasks via a single shared backbone has become a vital operational paradigm. To
route incoming queries to task-specific LoRA adapters on-the-fly, existing
dynamic model ensembling methods have grown increasingly convoluted—introducing
offline calibration datasets, classification-head dependencies, and parametric
density models like Gaussian Mixture Models (GMMs) for out-of-distribution (OOD)
rejection. In this paper, we advocate for a return to mathematical simplicity.
We present \textbf{LoRA Subspace Projection Routing (LSPR)}. Unlike post-hoc ensembling methods that attempt to serve arbitrary, standard pre-trained LoRA weights (which fail under LSPR due to the lack of weight-activation alignment), LSPR is a co-designed joint training-and-routing framework. It pairs fine-tuning with a lightweight representational autoencoding objective during training, which mathematically forces the column space of the down-projection weight matrix $A_k$ to span the task's activation subspace. By performing a microsecond-level, closed-form QR
decomposition of the early-layer LoRA down-projection matrices offline ($A_k =
Q_k R_k$), LSPR extracts an orthonormal basis $Q_k$ representing each task's
intrinsic representational subspace. At inference time, routing and OOD
rejection are resolved simultaneously by projecting early activations $h_b$ onto
these subspaces to yield scale-invariant geometric alignment scores. Evaluated
within a controlled, synthetic PyTorch multi-task environment (the Isolating
Coordinate Sandbox) designed as a proof-of-concept for representation geometry,
LSPR recovers 85.81\% Joint Mean accuracy (exceeding Uniform Merging by 61.85\% and perfectly recovering the expert ceiling and SPS-ZCA SOTA) and achieves an outstanding zero-shot OOD rejection AUROC of 0.9755, while delivering flat, highly efficient physical latency scaling on edge
CPUs. Our approach demonstrates that by pairing
fine-tuning with a lightweight representational autoencoding objective, elegant
linear algebra can resolve dynamic routing and OOD rejection at deployment time
with zero trainable parameters, zero classification-head dependencies, and zero task-specific calibration data.
\end{abstract}
